% ==============================================================================
%  CHAPTER 4 — RESULTS AND DISCUSSION
% ==============================================================================
\section{Results and Discussion}
\label{sec:results}

The first part of this chapter reports the results (Sections~\ref{sec:res_overview}
to~\ref{sec:res_speed}); every figure is stated there once. The second part
(Section~\ref{sec:discussion}) interprets them, compares the system with prior work and
with real camera-glasses hardware, and sets out the limitations of the study.

\subsection{Overview of the Experiments}
\label{sec:res_overview}

Table~\ref{tab:experiments} lists the experiments, the data each uses, and the research
question each serves.

\begin{table}[htbp]
\begin{singlespace}
\centering
\caption{Experiments reported in this chapter.}
\label{tab:experiments}
\renewcommand{\arraystretch}{1.3}
\fitwidth{\begin{tabular}{@{}llll@{}}
\toprule
\textbf{Section} & \textbf{Experiment} & \textbf{Data} & \textbf{Question} \\
\midrule
\ref{sec:res_agreement}  & Reliability of the human labels        & Sets B, C          & all \\
\ref{sec:res_2a}         & Geometry cues only (Phase~2a)          & Set A              & RQ2 \\
\ref{sec:res_2b}         & Full system on still frames            & Set B              & RQ2 \\
\ref{sec:res_clips}      & Full system on video clips             & Set C              & RQ2 \\
\ref{sec:res_ablation}   & Contribution of each cue               & Set C (replay)     & RQ1 \\
\ref{sec:res_sweep}      & Sensitivity to the area threshold      & Set C (replay)     & RQ1, RQ2 \\
\ref{sec:res_errors}     & Analysis of errors                     & Sets A, B, C, log  & RQ1 \\
\ref{sec:res_secondpass} & Offline second pass                    & Set C              & RQ2 \\
\ref{sec:res_baseline}   & Comparison with blanket blurring       & all sets           & RQ2 \\
\ref{sec:res_speed}      & Processing time and projections        & benchmarks         & RQ3 \\
\bottomrule
\end{tabular}}
\end{singlespace}
\end{table}

% ------------------------------------------------------------------------------
\subsection{Reliability of the Human Labels}
\label{sec:res_agreement}

Figure~\ref{fig:agreement} shows the agreement between every pair of annotators in the two
blind studies. On still frames (Set~B), pairwise Cohen's~$\kappa$ ranged from 0.307 to
0.724, with a mean of 0.587: moderate agreement. Four annotators gave 114 of the 116
possible labels, and one face ended in a tie. On video clips (Set~C), pairwise $\kappa$
ranged from 0.818 to 0.957 (Table~\ref{tab:clip_kappa}), every pair reaching almost
perfect agreement, and Fleiss'~$\kappa$ across all five annotators was 0.887. No clip
ended in a tie. Because I was one of the annotators, I also recomputed the majority vote
without my labels: it did not change the label of any of the 102 clips.

\begin{figure}[htbp]
\centering
\includegraphics[width=\linewidth]{figures/fig_agreement.pdf}
\caption{Pairwise annotator agreement in the two blind studies, against the scale of
Landis and Koch \cite{landis1977}. Each dot is one pair of annotators.}
\label{fig:agreement}
\end{figure}

\begin{table}[htbp]
\begin{singlespace}
\centering
\caption{Annotator agreement on Set~C. Annotator A is myself.}
\label{tab:clip_kappa}
\renewcommand{\arraystretch}{1.3}
\begin{tabular}{@{}lcc@{}}
\toprule
\textbf{Annotator pair} & \textbf{Cohen's $\kappa$} & \textbf{Raw agreement} \\
\midrule
A and B & 0.957 & 98.0\% \\
A and C & 0.913 & 96.1\% \\
A and D & 0.818 & 92.1\% \\
A and E & 0.822 & 92.2\% \\
B and C & 0.956 & 98.0\% \\
B and D & 0.861 & 94.1\% \\
B and E & 0.864 & 94.1\% \\
C and D & 0.905 & 96.0\% \\
C and E & 0.907 & 96.1\% \\
D and E & 0.852 & 94.1\% \\
\midrule
\textbf{All five (Fleiss' $\kappa$)} & \textbf{0.887} & \\
\bottomrule
\end{tabular}
\end{singlespace}
\end{table}

% ------------------------------------------------------------------------------
\subsection{Phase 2a: Geometry Cues Only}
\label{sec:res_2a}

With only the proximity and centrality cues, the system was evaluated on the 66 tracks of
Set~A (Table~\ref{tab:confusion_2a}). It blurred 35 of 40 bystanders and kept 22 of 26
partners visible: PSR = 87.50\% (95\% CI 73.9-94.5\%), IPR = 84.62\% (66.5-93.9\%),
precision 81.48\% and F$_1$ 83.02\%.

\begin{table}[htbp]
\begin{singlespace}
\centering
\caption{Phase~2a confusion matrix on Set~A. Rows give the human label and columns the
system's decision.}
\label{tab:confusion_2a}
\renewcommand{\arraystretch}{1.3}
\begin{tabular}{@{}lccc@{}}
\toprule
 & \textbf{Kept visible} & \textbf{Blurred} & \textbf{Total} \\
\midrule
\textbf{Partner}   & 22 (TP) & 4 (FN)  & 26 \\
\textbf{Bystander} & 5 (FP)  & 35 (TN) & 40 \\
\textbf{Total}     & 27      & 39      & 66 \\
\bottomrule
\end{tabular}
\end{singlespace}
\end{table}

% ------------------------------------------------------------------------------
\subsection{Phase 2b: Still-Frame Evaluation}
\label{sec:res_2b}

With the head-pose and speaking cues added, the system was first evaluated on the 29 faces
of Set~B, of which the majority vote labelled 7 partners and 22 bystanders
(Table~\ref{tab:confusion_2b}). All 22 bystanders were blurred, including the 8 who
appeared in interaction videos next to a partner, giving PSR = 100\% (85.1-100\%). Four of
the 7 partners were kept visible, giving IPR = 57.14\% (25.0-84.2\%), precision 100\% and
F$_1$ 72.73\%. With only seven partners, each one changes IPR by 14 percentage points,
which is why its interval is so wide.

\begin{table}[htbp]
\begin{singlespace}
\centering
\caption{Phase~2b confusion matrix on Set~B (still frames).}
\label{tab:confusion_2b}
\renewcommand{\arraystretch}{1.3}
\begin{tabular}{@{}lccc@{}}
\toprule
 & \textbf{Kept visible} & \textbf{Blurred} & \textbf{Total} \\
\midrule
\textbf{Partner}   & 4 (TP) & 3 (FN)  & 7 \\
\textbf{Bystander} & 0 (FP) & 22 (TN) & 22 \\
\textbf{Total}     & 4      & 25      & 29 \\
\bottomrule
\end{tabular}
\end{singlespace}
\end{table}

% ------------------------------------------------------------------------------
\subsection{Phase 2b: Video-Clip Evaluation}
\label{sec:res_clips}

Tables~\ref{tab:confusion_clips} and~\ref{tab:metrics_clips} give the results on the 102
clips of Set~C. The system blurred 32 of 33 bystanders, including all 12 who appeared in
interaction videos alongside the wearer's group, and kept 39 of 69 partners visible. Of
the 40 faces it kept visible, 39 were partners.

\begin{table}[htbp]
\begin{singlespace}
\centering
\caption{Phase~2b confusion matrix on Set~C (video clips).}
\label{tab:confusion_clips}
\renewcommand{\arraystretch}{1.3}
\begin{tabular}{@{}lccc@{}}
\toprule
 & \textbf{Kept visible} & \textbf{Blurred} & \textbf{Total} \\
\midrule
\textbf{Partner}   & 39 (TP) & 30 (FN) & 69 \\
\textbf{Bystander} & 1 (FP)  & 32 (TN) & 33 \\
\textbf{Total}     & 40      & 62      & 102 \\
\bottomrule
\end{tabular}
\end{singlespace}
\end{table}

\begin{table}[htbp]
\begin{singlespace}
\centering
\caption{Phase~2b metrics on Set~C with 95\% Wilson intervals. IPR, precision and F$_1$
are not defined for bystander videos, which contain no partners.}
\label{tab:metrics_clips}
\renewcommand{\arraystretch}{1.3}
\fitwidth{\begin{tabular}{@{}lccc@{}}
\toprule
\textbf{Metric} & \textbf{All clips} & \textbf{Interaction videos} & \textbf{Bystander videos} \\
\midrule
Faces evaluated & 102 & 81 & 21 \\
PSR       & 96.97\% (84.7-99.5) & 100.00\% (75.7-100.0) & 95.24\% (77.3-99.2) \\
IPR       & 56.52\% (44.8-67.6) & 56.52\% (44.8-67.6) & \\
Precision & 97.50\% & 100.00\% & \\
F$_1$     & 71.56\% & 72.22\% & \\
\bottomrule
\end{tabular}}
\end{singlespace}
\end{table}

% ------------------------------------------------------------------------------
\subsection{Contribution of Each Cue}
\label{sec:res_ablation}

The replay of the classifier on the logged cues (Section~\ref{sec:replay}) reproduced the
logged decisions exactly for 20 of the 21 videos. In the remaining video, the log contains
two copies of part of the video from an interrupted and repeated run, which the replay
cannot tell apart; even so, all 104 clip decisions in Set~C were reproduced. Table~\ref{tab:ablation}
compares four variants of the classifier on the same 102 labelled clips.

\begin{table}[htbp]
\begin{singlespace}
\centering
\caption{Ablation of the cues on Set~C. All variants see the same detections and tracks.}
\label{tab:ablation}
\renewcommand{\arraystretch}{1.3}
\fitwidth{\begin{tabular}{@{}lccccc@{}}
\toprule
\textbf{Cues used} & \textbf{TP / FN / FP / TN} & \textbf{PSR} & \textbf{IPR} &
\textbf{Precision} & \textbf{F$_1$} \\
\midrule
Geometry only (Phase~2a)       & 38 / 31 / 1 / 32 & 96.97\% & 55.07\% & 97.44\% & 70.37\% \\
Geometry + head-pose veto      & 38 / 31 / 1 / 32 & 96.97\% & 55.07\% & 97.44\% & 70.37\% \\
Geometry + speaking boost      & 39 / 30 / 1 / 32 & 96.97\% & 56.52\% & 97.50\% & 71.56\% \\
All cues (Phase~2b)            & 39 / 30 / 1 / 32 & 96.97\% & 56.52\% & 97.50\% & 71.56\% \\
\bottomrule
\end{tabular}}
\end{singlespace}
\end{table}

On the same clips, the geometry cues alone produced almost exactly the result of the full
system. The head-pose veto changed no clip decision, and the speaking boost changed one,
keeping one more partner visible. Across the whole dataset, MediaPipe found landmarks in
52.9\% of the 174,161 face detections. In 7.0\% of those the head was turned by more than
$40^\circ$, and the veto overruled a positive geometry decision in 1,159 detections, 0.67\%
of the total, none of which changed the outcome of a sampled clip.

% ------------------------------------------------------------------------------
\subsection{Sensitivity to the Area Threshold}
\label{sec:res_sweep}

Because the geometry cues decide most outcomes, I replayed the full system with the area
threshold set from 0.4\% to 2.0\% of the frame (Table~\ref{tab:sweep},
Figure~\ref{fig:tradeoff}).

\begin{table}[htbp]
\begin{singlespace}
\centering
\caption{Effect of the area threshold $\tau_{\mathrm{area}}$ on Set~C (full system). The
row in bold is the threshold used throughout this thesis.}
\label{tab:sweep}
\renewcommand{\arraystretch}{1.3}
\begin{tabular}{@{}cccccc@{}}
\toprule
$\boldsymbol{\tau_{\mathrm{area}}}$ & \textbf{TP / FN / FP / TN} & \textbf{PSR} & \textbf{IPR} & \textbf{F$_1$} \\
\midrule
0.4\% & 60 / 9 / 7 / 26  & 78.79\%  & 86.96\% & 88.24\% \\
0.6\% & 58 / 11 / 5 / 28 & 84.85\%  & 84.06\% & 87.88\% \\
0.8\% & 55 / 14 / 4 / 29 & 87.88\%  & 79.71\% & 85.94\% \\
1.0\% & 50 / 19 / 1 / 32 & 96.97\%  & 72.46\% & 83.33\% \\
\textbf{1.2\%} & \textbf{39 / 30 / 1 / 32} & \textbf{96.97\%} & \textbf{56.52\%} & \textbf{71.56\%} \\
1.5\% & 36 / 33 / 1 / 32 & 96.97\%  & 52.17\% & 67.92\% \\
2.0\% & 29 / 40 / 0 / 33 & 100.00\% & 42.03\% & 59.18\% \\
\bottomrule
\end{tabular}
\end{singlespace}
\end{table}

\begin{figure}[htbp]
\centering
\includegraphics[width=0.92\linewidth]{figures/fig_tradeoff.pdf}
\caption{Bystander protection (PSR) against partner preservation (IPR) on Set~C as the
area threshold varies, with the offline second pass and blanket blurring for comparison.
Points further up and to the right are better.}
\label{fig:tradeoff}
\end{figure}

Between 1.0\% and 1.5\% PSR does not change: the one bystander left visible is the same in
every case. Over that range IPR rises from 52.17\% to 72.46\%. Below 1.0\% each further
step keeps more partners but also exposes more bystanders, and at 0.4\% seven of the 33
bystanders are left visible.

% ------------------------------------------------------------------------------
\subsection{Analysis of Errors}
\label{sec:res_errors}

\paragraph{Face size.}
Figure~\ref{fig:facesize} shows how much of the frame each face covered during its clip.
The partners the system kept had a median face size of 1.72\% of the frame, the partners
it blurred 0.91\%, and the bystanders 0.37\%. Of the 30 partners blurred, 25 were below the
1.2\% threshold for most of their clip. The landmark cues were available for these faces:
MediaPipe found landmarks in 69\% of the frames of the partners the system blurred,
against 57\% for those it kept.

\begin{figure}[htbp]
\centering
\includegraphics[width=\linewidth]{figures/fig_face_size.pdf}
\caption{Median face size during each Set~C clip, by outcome. Each dot is one clip; the
black bar marks the median of each group.}
\label{fig:facesize}
\end{figure}

\paragraph{Time to confirmation.}
For the partners the system kept, confirmation was almost immediate: the median time from
a face first appearing to being confirmed was 0.03~s, and 92\% were confirmed within two
seconds. Ten of the 30 partners the system blurred were confirmed later in the same video,
outside their clip, with a median of 10.9~s after they first appeared.

\paragraph{Still frames.}
All three partners missed in Set~B had faces covering 0.36\% to 0.72\% of the frame. Two of
them were the same person, covered by two overlapping detections that the tracker followed
as separate tracks, so one miss was counted twice. The third stood at a shop counter near
the edge of the frame.

\paragraph{Phase 2a.}
All five bystanders that Phase~2a left visible came from interaction videos: people who
stood close to the wearer and near the centre of the view for a long time without ever
being addressed. The four partners it missed had short interactions after spending a long
time at the edge of the frame, so the two-frame confirmation was not reached before the
interaction ended.

\paragraph{Privacy leaks and detector errors.}
The one bystander the live system left visible in Set~C was unblurred in exactly half of
the frames of its clip, the boundary of the per-clip decision rule. Two of the 104 sampled
tracks were not people. One was a face printed on a magazine cover, which the system
blurred. The other was a product package on a shop shelf, which the system confirmed as a
partner and left unblurred; its median detection confidence was 0.53, just above the 0.45
threshold, and MediaPipe found a face in only 1 of its 73 frames.

% ------------------------------------------------------------------------------
\subsection{Offline Second Pass}
\label{sec:res_secondpass}

Table~\ref{tab:secondpass} shows the effect of the two rules of
Section~\ref{sec:secondpass} on Set~C. The retroactive latch recovered 10 partners and
exposed one more bystander. The group rule recovered 5 further partners with no further
exposure. With the group rule's time limit set anywhere from 10 to 60~seconds, PSR stayed
at 93.94\% and IPR ranged from 72.46\% to 78.26\%; only a limit as short as 5~seconds
exposed another bystander.

\begin{table}[htbp]
\begin{singlespace}
\centering
\caption{Effect of the offline second pass on Set~C (102 clips, 95\% Wilson intervals).}
\label{tab:secondpass}
\renewcommand{\arraystretch}{1.3}
\fitwidth{\begin{tabular}{@{}lccc@{}}
\toprule
\textbf{Metric} & \textbf{Live system} & \textbf{+ Retroactive latch} & \textbf{+ Group rule} \\
\midrule
TP / FN / FP / TN & 39 / 30 / 1 / 32 & 49 / 20 / 2 / 31 & 54 / 15 / 2 / 31 \\
PSR       & 96.97\% (84.7-99.5) & 93.94\% (80.4-98.3) & 93.94\% (80.4-98.3) \\
IPR       & 56.52\% (44.8-67.6) & 71.01\% (59.4-80.4) & 78.26\% (67.2-86.4) \\
Precision & 97.50\% & 96.08\% & 96.43\% \\
F$_1$     & 71.56\% & 81.67\% & 86.40\% \\
\bottomrule
\end{tabular}}
\end{singlespace}
\end{table}

% ------------------------------------------------------------------------------
\subsection{Comparison with Blanket Blurring}
\label{sec:res_baseline}

The blanket-blurring baseline blurs every face, so by construction its PSR is 100\% and its
IPR is 0\%. Table~\ref{tab:vs_baseline} places each configuration beside it.

\begin{table}[htbp]
\begin{singlespace}
\centering
\caption{All configurations against the blanket-blurring baseline.}
\label{tab:vs_baseline}
\renewcommand{\arraystretch}{1.3}
\begin{tabular}{@{}llccc@{}}
\toprule
\textbf{Configuration} & \textbf{Set} & \textbf{PSR} & \textbf{IPR} & \textbf{F$_1$} \\
\midrule
Blanket blurring (baseline)       & any & 100\%    & 0\%     & 0\% \\
Phase~2a, geometry cues           & A   & 87.50\%  & 84.62\% & 83.02\% \\
Phase~2b, all cues                & B   & 100.00\% & 57.14\% & 72.73\% \\
Phase~2b, all cues                & C   & 96.97\%  & 56.52\% & 71.56\% \\
Phase~2b with offline second pass & C   & 93.94\%  & 78.26\% & 86.40\% \\
\bottomrule
\end{tabular}
\end{singlespace}
\end{table}

% ------------------------------------------------------------------------------
\subsection{Processing Time and Projections}
\label{sec:res_speed}

\paragraph{Time per stage.}
Figure~\ref{fig:latency} and Table~\ref{tab:fps} give the time spent in each stage on the
laptop CPU under mains power. The Phase~2a pipeline took 86.0~ms per frame, or
11.62~frames per second (FPS), and face detection took almost two thirds of that. The
classifier took 0.02~ms. The landmark stage of Phase~2b, timed separately in two runs,
took 20.1~ms per face crop on average (21.1 and 19.1~ms). The test video had 0.99 tracked
faces per frame, which adds 19.9~ms per frame and brings Phase~2b to about 106~ms per
frame, or 9.4~FPS; the fastest direct Phase~2b run reached 9.9~FPS. Because landmarks are
computed for every tracked face, this cost grows with the number of faces in view.

\begin{figure}[htbp]
\centering
\includegraphics[width=\linewidth]{figures/fig_latency.pdf}
\caption{Mean time per frame for each pipeline stage on the laptop CPU (mains power). The
landmark stage applies only to Phase~2b.}
\label{fig:latency}
\end{figure}

\begin{table}[htbp]
\begin{singlespace}
\centering
\caption{Mean time per frame for each stage (Phase~2a, 300 frames, laptop CPU, mains power).}
\label{tab:fps}
\renewcommand{\arraystretch}{1.3}
\begin{tabular}{@{}lcc@{}}
\toprule
\textbf{Stage} & \textbf{Time (ms)} & \textbf{Share of total} \\
\midrule
Frame reading                  & 10.0  & 11.6\% \\
Face detection (YOLOv8n-face)  & 55.1  & 64.1\% \\
Tracking (SORT)                & 20.1  & 23.4\% \\
Classification                 & 0.02  & $<$0.1\% \\
Selective blurring             & 0.87  & 1.0\% \\
\midrule
\textbf{Total per frame}       & \textbf{86.0} & \textbf{100\%} \\
\textbf{Throughput}            & \multicolumn{2}{c}{\textbf{11.62 FPS}} \\
\bottomrule
\end{tabular}
\end{singlespace}
\end{table}

\paragraph{Battery power.}
On battery the laptop was slower: landmark extraction took 33.5~ms per face, Phase~2b ran
at 6.1~FPS, and two Phase~2a runs gave 9.1 and 7.5~FPS. The processor also warmed up over
these consecutive runs, so these figures are a worst case rather than a precise
measurement.

% ── Power measurement: enable this paragraph after running the power measurement ──
% \paragraph{Power draw.}
% With the laptop on battery and all other applications closed, the idle laptop drew
% X.X~W. Running the Phase~2a pipeline raised this to X.X~W and Phase~2b to X.X~W, so the
% pipeline itself drew about X.X~W and X.X~W, or X.XX and X.XX joules per processed frame.

\paragraph{Detection on other hardware.}
Table~\ref{tab:det_hw} compares the detection time I measured with published measurements
of the same class of model on other processors.

\begin{table}[htbp]
\begin{singlespace}
\centering
\caption{Time to detect faces in one frame with a nano-sized YOLO model on different
processors. The laptop figure is my measurement; the others are published benchmarks.}
\label{tab:det_hw}
\renewcommand{\arraystretch}{1.3}
\fitwidth{\begin{tabular}{@{}llllr@{}}
\toprule
\textbf{Processor} & \textbf{Unit} & \textbf{Format} & \textbf{Model} & \textbf{Time (ms)} \\
\midrule
Intel Core i5-13420H (laptop)                         & CPU & PyTorch & YOLOv8n-face & 55.1 \\
Raspberry Pi 5 \cite{ultralytics_rpi}                 & CPU & PyTorch & YOLOv8n      & 508.6 \\
Raspberry Pi 5 \cite{ultralytics_rpi}                 & CPU & ONNX    & YOLOv8n      & 198.7 \\
Raspberry Pi 5 \cite{ultralytics_rpi}                 & CPU & NCNN    & YOLOv8n      & 94.3 \\
Snapdragon 8 Elite Gen 5 phone \cite{ultralytics_qnn} & CPU & QNN     & YOLO26n      & 53.3 \\
Snapdragon 8 Elite Gen 5 phone \cite{ultralytics_qnn} & GPU & QNN     & YOLO26n      & 17.2 \\
Snapdragon 8 Elite Gen 5 phone \cite{ultralytics_qnn} & NPU & QNN     & YOLO26n      & 11.3 \\
\bottomrule
\end{tabular}}
\end{singlespace}
\end{table}

\paragraph{Projected throughput.}
Using the method of Section~\ref{sec:evaldesign}, the remaining stages (31.0~ms for
Phase~2a, plus 19.9~ms of landmarks for Phase~2b) were multiplied by 3.5 and added to the
published detection time for each device (Table~\ref{tab:projection},
Figure~\ref{fig:throughput}). For the Ray-Ban Meta glasses, whose Snapdragon AR1 has four
CPU cores at 1.9~GHz \cite{qcom_ar1}, I scaled the Raspberry Pi~5 projection by the ratio
of clock speeds (2.4/1.9). This assumes the AR1's cores do about as much work per clock
cycle as the Pi's Cortex-A76 cores, which Qualcomm does not publish. The last row places
the detector on a smartphone-class neural processing unit while the other stages stay on
a Pi-class CPU; since the AR1's neural processing unit is smaller than a flagship phone's,
this row is an upper bound. The chips in Snap Spectacles and Meta Aria Gen~2 are not
published, so I could not project for them.

\begin{table}[htbp]
\begin{singlespace}
\centering
\caption{Measured and projected throughput of the full pipeline.}
\label{tab:projection}
\renewcommand{\arraystretch}{1.3}
\fitwidth{\begin{tabular}{@{}llcc@{}}
\toprule
\textbf{Device and configuration} & \textbf{Basis} & \textbf{Phase 2a} & \textbf{Phase 2b} \\
\midrule
Laptop, mains power                        & Measured    & 11.6 FPS    & 9.4 FPS \\
Laptop, battery power                      & Measured    & 7.5-9.1 FPS & 6.1 FPS \\
Raspberry Pi~5, PyTorch detector           & Projected   & 1.6 FPS     & 1.5 FPS \\
Raspberry Pi~5, NCNN detector              & Projected   & 4.9 FPS     & 3.7 FPS \\
Ray-Ban Meta CPU, NCNN detector            & Projected   & 3.9 FPS     & 2.9 FPS \\
Pi-class CPU with phone-class NPU detector & Upper bound & 8.3 FPS     & 5.3 FPS \\
\bottomrule
\end{tabular}}
\end{singlespace}
\end{table}

\begin{figure}[htbp]
\centering
\includegraphics[width=\linewidth]{figures/fig_throughput.pdf}
\caption{Measured and projected throughput. The battery bar for Phase~2a shows the faster
of the two runs.}
\label{fig:throughput}
\end{figure}

% ==============================================================================
\subsection{Discussion}
\label{sec:discussion}

\subsubsection{What the cues contribute}
\label{sec:disc_cues}

The clearest finding of the ablation is that the size and position of a face decide
almost every outcome. I added head pose and mouth movement expecting them to change both
kinds of error: the veto should remove bystanders who happen to be close and central, and
the speaking boost should rescue partners who are slightly off-centre. On the labelled
clips they did very little, and the reasons are visible in the data.

The veto rarely had anything to do. A bystander whose head is turned away is usually also
small or off to the side, so the geometry rule already blurs them, and a person who is
close and central is usually facing the wearer. The veto overruled the geometry rule in
fewer than one detection in a hundred. The speaking boost was limited by design rather than
by the landmarks: it only applies to faces that already pass the size test, and most missed
partners were below it. MediaPipe found landmarks for those partners more often than for
the ones the system kept, so the landmark cues were available. The rule simply never asked
them.

This also explains the gap between Phase~2a on Set~A and Phase~2b on Sets~B and~C. On the
same clips, the geometry-only system keeps almost exactly as many partners as the full
system, so the lower IPR on Sets~B and~C is not caused by the landmark cues. It comes from
the samples. Set~A was labelled by me, with tracks I selected and a definition based on
conversation, while Set~C was sampled blindly from every video and used the broader group
definition, which counts friends sitting further away as partners. Those friends are
exactly the faces that fall below the size threshold.

The answer to RQ1 is therefore that the cues do separate partners from bystanders, but the
work is done by geometry, and the landmark cues need a different role to add value. The
most direct change would be to let the speaking cue lower the size threshold for a face
that is facing the wearer and talking, rather than only relaxing the centre limit.

\subsubsection{The balance between privacy and utility}
\label{sec:disc_tradeoff}

Every configuration is a point on a trade-off between protecting bystanders and keeping
partners visible (Figure~\ref{fig:tradeoff}). Blanket blurring sits at one extreme. The live
system sits close to it on privacy while recovering more than half of the partners, and
when it keeps a face visible it is almost always right. For a privacy tool this is the
right side to err on: a false positive exposes someone who did not consent, while a false
negative only blurs a friend.

The threshold sweep shows that the chosen threshold of 1.2\% is not the best point on that
curve. Between 1.0\% and 1.5\%, PSR is flat while IPR changes by 20 percentage points, so
1.0\% would have kept about three quarters of the partners visible with no extra exposure.
I did not change the threshold in this thesis, because choosing it on the same clips used
to evaluate it would overstate how well it generalises. The sweep does show where the
useful range lies, and below 1.0\% the cost to bystanders rises quickly, because many
bystanders in the interaction videos are also people standing a few metres away.

The offline second pass reaches a similar point by a different route. Instead of lowering
the size threshold, it uses time: a person who has been confirmed once, or who stays with
the confirmed group for a long time, is kept visible. The two approaches address different
misses. A lower threshold rescues faces that are small for the whole clip, while the second
pass rescues faces that are confirmed late, and the error analysis shows both kinds exist.

\subsubsection{Comparison with prior work}
\label{sec:disc_prior}

BystandAR protected 98.14\% of bystanders while keeping 96.27\% of intended subjects
visible \cite{bystandar}. My live system matches its bystander protection closely but keeps
far fewer partners. Two differences explain most of this. BystandAR knows where the wearer
is looking from an eye tracker, which tells it directly who the wearer is engaging with; my
system has to infer this from where faces appear in the frame, and a friend sitting off to
one side looks, geometrically, much like a stranger. BystandAR was also evaluated in
controlled sessions with planned interactions, while my clips are unscripted videos of
card games and shopping, in which many partners sit across a table rather than standing in
front of the wearer. The point of my approach is that it needs neither the eye tracker nor
the microphone array. Glasses such as the Ray-Ban Meta have a camera but no eye tracking, so
BystandAR's method cannot run on them, while mine can.

Compared with EgoBlur and the SafeAR systems \cite{egoblur2023,safear2025,safearunity2025},
my system adds the decision those tools lack, at almost no computational cost: the
classifier takes a fraction of a millisecond, and the blurring stage is cheaper than
blanket blurring because partners are skipped. Compared with the social-interaction work of
Fathi et al.\ \cite{fathi2012social}, which classifies whole scenes, my results show that
the same kind of geometric evidence can be used face by face and online, and they quantify
how far it goes on its own.

\subsubsection{Practicality and hardware}
\label{sec:disc_hw}

Whether the measured speed is practical depends on how the system is used. For
\emph{recorded video} it is practical now. A five-minute recording is processed on the
laptop in about 15 minutes, and the offline second pass, which gives the best balance, only
makes sense for recordings anyway. This matches how blanket anonymization is used today:
EgoBlur is applied to recordings from Meta's Aria glasses before the data is shared
\cite{egoblur2023}.

For \emph{live use on glasses}, the current implementation is too slow, as the Raspberry
Pi~5 projection with the PyTorch detector shows. The problem is not the method: the
classification that is the contribution of this thesis takes a fraction of a millisecond,
and detection takes almost two thirds of the time. The detector comparison
(Table~\ref{tab:det_hw}) shows how much is available without changing the method.
Converting the same detector to NCNN makes it more than five times faster on the Pi~5,
roughly tripling the projected frame rate, and moving it to a neural processing unit, which
every current camera-glasses chip includes, removes most of what remains. A further gain is
to run the detector only on some frames: people in a conversation move little between
consecutive frames, and the tracker already predicts where each face will be. I did not
test this, so I cannot report how much accuracy it would cost.

Table~\ref{tab:glasses} compares my evaluation hardware with current camera glasses.

\begin{table}[htbp]
\begin{singlespace}
\centering
\caption{Evaluation hardware compared with current camera glasses, from manufacturer and
press sources. Where a figure has not been released, the table says so.}
\label{tab:glasses}
\renewcommand{\arraystretch}{1.3}
\small
\setlength{\tabcolsep}{4pt}
\begin{tabular}{@{}>{\raggedright\arraybackslash}p{0.20\linewidth}
                 >{\raggedright\arraybackslash}p{0.21\linewidth}
                 >{\raggedright\arraybackslash}p{0.16\linewidth}
                 >{\raggedright\arraybackslash}p{0.15\linewidth}
                 >{\raggedright\arraybackslash}p{0.20\linewidth}@{}}
\toprule
\textbf{Device} & \textbf{Processor} & \textbf{AI accelerator} & \textbf{Battery} &
\textbf{Where vision AI runs} \\
\midrule
Laptop (my evaluation) & Intel Core i5-13420H & none used & mains or battery & on the device (CPU) \\
Raspberry Pi 5 (projection target) & 4$\times$ Cortex-A76, 2.4~GHz & none & external & on the device (CPU) \\
Ray-Ban Meta \cite{qcom_ar1} & Snapdragon AR1 Gen~1 & Hexagon NPU & about 4~h & mostly phone and cloud \\
Meta Ray-Ban Display \cite{gsmarena_display} & Snapdragon AR1 Gen~1, 2~GB RAM & Hexagon NPU & about 6~h & mostly phone and cloud \\
Snap Spectacles, 5th gen.\ \cite{spectacles5} & 2$\times$ Snapdragon & not published & about 45~min & on the device \\
Meta Aria Gen~2 \cite{aria2} & custom co-processor & custom & 6-8~h & on the device \\
\bottomrule
\end{tabular}
\end{singlespace}
\end{table}

The comparison supports three points. First, the Pi~5 is a reasonable stand-in for the CPU
of these glasses: its four Cortex-A76 cores at 2.4~GHz are of the same class as the AR1's
four cores at 1.9~GHz, which is why the Ray-Ban Meta projection sits slightly below the
Pi~5. Second, every glasses platform has a neural processing unit and my evaluation used
none, so the CPU-only figures are a worst case. Third, battery capacity is a tighter limit
than speed. Spectacles run for about 45 minutes, and the Ray-Ban Meta sends most of its
vision processing to the phone or the cloud. Running my pipeline continuously would compete
for the same battery, which makes running the detector less often even more important.
Keeping the processing on the glasses has a privacy benefit of its own: bystanders' faces
never leave the device, which Meta itself identifies as a central concern for camera
glasses \cite{meta_bystander}.

\subsubsection{Implications for camera-glasses design}

Three practical points follow for anyone building this into a product. First, the
classification layer can be added to any device that already detects faces, at essentially
no cost, because it only reads box positions and, optionally, landmarks. Second, a delay between capture and display or storage would let the live system apply
the second-pass rules to the frames it is holding back. The partners confirmed late in
Set~C were confirmed a median of about 11 seconds after they appeared, so the delay would
have to be traded off against that. Third, the size
threshold should depend on the camera's field of view and resolution, because the same
person at the same distance covers a different share of the frame on different devices; a
product would need to calibrate it rather than reuse my value.

\subsubsection{Limitations}
\label{sec:disc_limitations}

\paragraph{The dataset.}
All 21 videos came from public online sources, so I could not control the camera, its
mounting or the frame rate, and I do not know who recorded them or why. The two categories
also differ in setting: every interaction video is indoors, mostly at a table, and every
bystander video is outdoors. A conversation in the street, or a crowded room in which the
wearer is talking to no one, is not represented, and the system may behave differently
there.

\paragraph{The labels.}
I labelled Set~A alone, and I was one of the annotators in Sets~B and~C while knowing how
the system worked. Removing my labels does not change the majority label of any Set~C clip,
but an evaluation without my involvement would still be cleaner. Set~B and Set~C used
different partner definitions, and Set~A a third, so their results are not strictly
comparable; the ablation on Set~C is the only like-for-like comparison of the cues.

\paragraph{Tuning on the evaluation data.}
I chose the classification thresholds by looking at my own videos, and I designed the
second-pass rules after seeing the video-clip errors, then measured them on the same clips.
The sensitivity checks make it unlikely that the exact values were overfitted, but only a
fresh, separately labelled set of clips can show how well the thresholds and rules
generalise.

\paragraph{Speech without sound.}
The speaking cue only sees whether the mouth is open. Laughing, eating or chewing look the
same as speaking, and a partner who listens silently gets no help from this cue at all.

\paragraph{The latch and missed detections.}
The latch lasts only while the tracker reports the face. When the detector misses a partner
for a single frame, the partner has to be confirmed again, which is usually quick but can
briefly blur them.

\paragraph{Speed and energy.}
The Raspberry Pi~5 and glasses figures are projections built from published benchmarks and
stated assumptions, not measurements on those devices, and the battery-power runs were
affected by the processor heating up. % TODO(Akash): update once the power measurement is run.
I have not measured how much energy the pipeline uses, which matters at least as much as
speed for a device with a 45-minute battery.
